\SourceLeaf{336}%
\textsc{Bessel} for \(a\) in toises (Art. 5), add the following to the logarithm if \(r\) was originally expressed in toises
\[
0{,}2897827662
\]
or, if \(r\) was originally given in French legal meters, subtract the following from the logarithm
\[
0{,}0000371638
\]

The terms containing the factors (3), (4), and so on can be regarded as corrections that turn the approximate logarithms \(\log\tau,\log\beta,\log\lambda\) into the corrected logarithms \(\log t,\log b,\log l\). In all cases for which our method is intended, these corrections are very small decimal fractions. Since the logarithms are usually computed to seven figures, it is convenient to obtain the corrections directly in units of the seventh decimal place. To do this, one multiplies the values assigned to coefficients (3), (4), and so on in the preceding article by ten million, or increases their logarithms by seven units. This is how they are entered in our table. Accordingly, we also set
\[
\begin{aligned}
\log(7)&=4{,}62872\ (-10)\\
\log3(7)&=5{,}10584\ (-10)
\end{aligned}
\]
Even with this change, (3), (4), (5), (6), like (1) and (2), are still proper fractions, so their logarithms are themselves negative. They appear in the table in the usual way, with ten units borrowed for every logarithm.

\Article{35}
Everything said above (Art. 20) about solving the same problem on the sphere applies almost unchanged, with only minor adjustments, to the use of our formulas to solve the problem posed at the beginning of this memoir. Denote the quantities actually given, namely the latitude and azimuth at the first point, by \(B^0\) and \(T^0\). First, starting from an approximate value of \(T\) (if no other information is available, %
